Following precedent in employment discrimination law,
two notions of disparity are widely-discussed 
in papers on fairness and ML. 
Algorithms exhibit \emph{treatment disparity}
if they formally treat 
members of protected subgroups differently;
algorithms exhibit \emph{impact disparity}
when outcomes differ across subgroups (even unintentionally).
Naturally, we can achieve impact parity 
through purposeful treatment disparity.
One line of papers aims to reconcile the two parities
proposing \emph{disparate learning processes} (DLPs).
Here, the sensitive feature is used during training 
but a \emph{group-blind} classifier is produced.
In this paper, we show that:
(i) when sensitive and (nominally) nonsensitive features are correlated,
DLPs will indirectly implement treatment disparity,
undermining the policy desiderata 
they are designed to address;
(ii) when group membership is \emph{partly} revealed by other features, DLPs induce within-class discrimination;
and (iii) in general, DLPs provide suboptimal trade-offs between accuracy and impact parity.
Experimental results on several real-world datasets 
highlight the practical consequences 
of applying DLPs. 

